\section{Problem statement}
The study conducted by Srinivasa Ragavan suggests, that information detouring contributes to approximately 40\% of the time to understand a Spreadsheet. In some cases participants even failed to obtain the relevant information. Information seeking involves identifying dependants of a cell, understanding the underyling logic of a certian value. Another time consuming task is formula traccing. This occurs when a formula refers to another cell address, which again refers another and by doing so creating a chain of dependencies. Although modern spreadsheet programs highlight referenced cells on the grid, it is not sufficient when the cells are not visible in the grid\cite{srinivasa_ragavan_spreadsheet_2021}. 
A second problem is the loss of context which occurs during the transition from the author to the reader. An author has implicit knowledge about provenance and relation of the data which rarely gets encoded in the spreadsheet. Consequently, the reader lacks the context to fully reconstruct the relationships and provenances \cite{kohlhase_context_nodate}.
These challenges highlight a gap in the current spreadsheet programs. While there is a clear need to improve the visualization for the logic of a Spreadsheet, visualization alone is not enough. To address the problem of information detouring a tool must actively support the user in the navigation process, to make it easier to find relevant information while minimizing the context loss of the process.



%comprehension process is the information detouring. Where the dependencies and logic of the related data is not directly visible to the user \cite{srinivasa_ragavan_spreadsheet_2021}.
%Another problem for comprehending a Spreadsheet is the lost context when going from author to reader. The author has more context on the provenance on the data of the Spreadsheet than the reader which implies that there is a lack of presenting the locality and relationship of the data \cite{kohlhase_context_nodate}. These problems indicate first that there is still a need for a tool to improve the presenting and visualizing of the underlying logic of a Spreadsheet. But visualizing alone is not enough. The tool also has to support the navigation process without hindering to build a mental model of the system.